% chap1.tex
\chapter{序論}

\section{研究背景 - 水系感染症の現状 -}

安全な飲料水や衛生設備が十分に整っていない地域では、細菌による水系感染症が多発している。
そのため、それらの地域では安全で清潔な水へのアクセスが非常に困難であり、
健康被害が深刻化している\textsuperscript{1)}
一方で、日本は水道の普及率が98.3\%と非常に高く\textsuperscript{2)}、
ほとんどの地域で清潔で安全な水を安定して利用することができる。
この点からも、日本は他国と比較して水資源の確保や衛生面で大きく進んでいることがわかる。

しかし、世界全体で見ると、いまだに多くの人々が安全で清潔な水を十分に利用できない状況にある。
汚染された水を飲用あるいは生活用水として使用することで感染症にかかる人が後を絶たず、
特に発展途上国では深刻な社会問題となっている。
代表的な水系感染症としては、コレラ、赤痢、チフス、A型肝炎などが挙げられ、
これらはいずれも汚れた水や不衛生な環境が主な原因で発生する。
世界保健機関（WHO）の報告\textsuperscript{3)}によると、
安全な飲料水を利用できない人は世界で約20億人以上にのぼり、
毎年、約505,000人が安全でない水および不十分な衛生環境を原因として死亡している。
従って、安全な水の確保は依然として世界的な課題である。

また、安全な水資源へのアクセスが困難となる要因として、
多くの発展途上国では水道管や浄水場といったインフラが十分に整備されていないことが挙げられる。
さらに、設備の老朽化や下水・工場排水の河川への直接流入によって水源が汚染される事例も多く、
こうした汚染は水系感染症の発生要因となり、健康面でも重大な影響を及ぼしている。

加えて、気候変動の影響による干ばつの増加や地下水の過剰利用によって、
利用可能な水資源そのものが減少している地域も存在する。
貧困により水道料金を支払えない、あるいは浄水設備を利用できない家庭も多く、
結果として不衛生な水に依存せざるを得ない状況が生じている。
紛争や政治的不安定によるインフラ破壊や資源管理の不全が重なる地域では、
状況はさらに深刻化する。

具体例として、2010年のハイチ大地震後には下水処理機能の喪失により汚染水が広がり、
大規模なコレラ流行が発生した。
また、内戦下のイエメンでは水道インフラの破壊と医療体制の崩壊により、
世界最大規模のコレラ流行が確認されている。
南アジアのスラム地域では上下水道未整備による飲料水汚染から赤痢が季節的に多発し、
難民キャンプにおいてもトイレ不足や水の共同利用による集団感染が報告されている。
これらはいずれも、災害・紛争・貧困・インフラ不足による衛生環境の悪化が
感染拡大の主要因となっている事例である。

日本においても、東日本大震災の際には水道が停止し、
手洗い用の水が不足したことで十分な手指衛生が困難となり、
ノロウイルスや感染性胃腸炎の集団発生が報告された。
このように、平時には安全とされる地域であっても、
災害時には水系感染症のリスクが急激に高まることが示されている。

%--------------------------------------------------

\section{既存研究の限界}

現在までに発表されている研究の多くは、
日常生活環境や医療現場など、比較的衛生状態が管理された状況を対象としており、
災害時や水道が利用できない状況、高度に汚染された水環境を想定した実験データは
十分に蓄積されていない。

人が日常生活の中でさまざまな物体に手で触れた際に、
どの程度の微生物が皮膚表面に付着するかを測定した研究はいくつか存在するが、
環境微生物の種類や濃度が大きく変化する災害時の状況に
そのまま適用することは困難である。
そのため、災害時を想定した実験的検討が求められている。

本節では、これまでに報告されている代表的な既存研究について概説する。

\paragraph{病院エレベーターの押しボタンと手指を対象とした菌の伝播に関する考察}\textsuperscript{4)}

松野 容子、水野 秀一、大徳 優子による研究では、
病院環境における菌の分布および動態を明らかにすることを目的として、
病棟エレベーターの押しボタンと手指を対象とした細菌学的検討が行われた。
使用人数別（1～8人）および経時間的（30～120分）に押しボタン上から検出された菌数は、
概ね0～300 CFUであり、その70\%以上が20 CFU未満であった。
一方で、最大1200 CFUが検出される事例も確認され、
一部の手指汚染状況によっては極めて高い汚染が生じ得ることが示された。

\paragraph{高度汚染した手指の衛生学的手洗いの検討}\textsuperscript{5)}

矢野 久子、小林 寛伊、奥住 捷子は、
高度に手指が汚染された状況を想定し、
手洗い時間や手洗い方法の違いによる除菌効果について検討を行った。
本研究では、病院現場における手洗い頻度や時間の短さが課題であることが指摘されている。

\paragraph{手洗い効果の細菌学的考察}\textsuperscript{6)}

石田 和夫、三浦 英雄は、
1997～1999年に学生399人を対象として、
水洗い、石鹸、消毒薬など複数の方法による手洗い前後の
手指生菌数を比較し、手洗い方法による除菌効果の違いを明らかにした。

\paragraph{高齢者介護施設における入居者の手指衛生の実態}\textsuperscript{7)}

中村 明世、岡田 淳子、飯田 忠行は、
高齢者介護施設における入居者の手指衛生の実態を調査し、
活動状況に応じて手指衛生方法や実施頻度が異なることを報告している。

%--------------------------------------------------

\section{本研究の目的と重要性}

災害発生時や安全な水資源へのアクセスが制限された状況では、
衛生環境が急激に悪化し、感染症リスクが著しく増大する。
しかし、このような特殊な状況を対象とした微生物学的研究は、
現時点では十分に行われていない。

特に災害時には、手洗いに使用できる清潔な水の確保が困難となるため、
通常時と比較して手指衛生の効果が低下する可能性がある。
また、現在の防災マニュアルにおいても、
使用する水の種類や量によってどの程度微生物が除去されるのかといった
科学的根拠は十分に示されていない。

そこで本研究では、災害時に水道などの安全な水源の利用が困難となり、
河川水、雨水、貯留水などの代替水源を使用する状況を想定し、
水を媒介とした水系感染症のリスク評価を行うことを目的とする。
具体的には、以下の事項について検討を行った。

\begin{itemize}
  \item 模擬水を用いた掌への水分付着量の定量化
  \item 大腸菌ファージを指標微生物として用いた掌への移行率評価
  \item 手洗い後の掌への微生物残存率の評価
  \item 統計解析および定量的微生物リスク評価（QMRA）の実施
\end{itemize}

得られた結果をもとに、災害時における感染症発生リスクを評価し、
安全な水資源へのアクセスが制限される状況においても実践可能な
効果的な手指衛生方法および感染症予防対策を提案することを
本研究の最終的な目的とする。

%--------------------------------------------------

\section{論文構成}

本論文は以下の構成からなる。

本研究は、災害時など安全な水資源が十分に確保できない状況において、
手指を介した水系感染症のリスクを定量的に評価することを目的として
実施された。\\

まず第1章では、水系感染症の現状と手指衛生の重要性について整理し、
特に災害時や不衛生な水環境下において、手指への付着量や洗浄後の残存量が十分に定量評価されていないという
課題を明らかにし、本研究の目的と意義を示した。\\ 

次に第2章では、以降の実験結果を被験者間で適切に比較するための前提条件として、
掌面積および掌体積の評価手法を確立した。具体的には、スキャナー画像と
画像処理ソフトを用いて掌面積を算出し、さらに掌厚の測定による体積推定を行った。
また、測定結果の再現性や左右差について検討し、後続の章における付着量および
残存量の正規化と比較解析の基盤を整えた。\\

第3章では、掌が水に接触した際に皮膚表面に保持される水分量を定量的に測定した。
右手および左手について繰り返し測定を行い、被験者間差や左右差の有無を統計的に解析した。
さらに、第2章で求めた掌面積を用いて水分付着量の正規化を行い、掌の大きさに依存しない付着特性の評価を試みた。\\

第4章では、指標微生物として大腸菌ファージを用い、洗浄条件の違いがウイルス除去効率に及ぼす影響を検討した。
特に、分割洗浄と一括洗浄の比較を通じて、洗浄方法および使用水量の違いがウイルス除去挙動に与える
影響を明らかにした。\\

続く第5章では、ウイルス付着実験および洗浄実験で得られたデータを整理し、
掌面積による標準化を行った上で被験者間比較を実施した。
また、洗浄回数の増加に伴うウイルス残存率の変化を解析し、ウイルス除去挙動の傾向を明確にした。\\ 

第6章では、第3章および第5章で得られた実測データを基に、
定量的微生物リスク評価（QMRA）を実施した。曝露シナリオを設定し、
モンテカルロシミュレーションを用いて年間感染確率を推定することで、
洗浄後に掌へ残存するウイルス量が感染リスクに与える影響を評価した。\\

第7章では、本研究全体の結果を総合的に整理し、得られた知見の意義および限界について考察した。
さらに、災害時や水資源が制限される状況における手指衛生対策への示唆を示すとともに、今後の課題について述べた。 
